\documentclass[12pt]{article}
\usepackage[utf8]{inputenc}
\usepackage[T1]{fontenc}
\usepackage{hyperref}
\title{Western Tarahumara Language Resources for K-12 Curriculum Development}
\author{Compiled by Claude Sonnet 5 (claude-sonnet-5) for Andr\'as Kornai}
\date{}
\begin{document}
\maketitle

\section{Introduction}
Western Tarahumara (ISO 639-3 \texttt{tac}, Glottocode lowl1265) is a Uto-Aztecan language of Mexico with about 10,000 speakers. It is not known to be taught in schools. The search found Scripture and audio resources. It found no dictionary or school materials. Do not confuse this variety with the Samachique (Central) Tarahumara New Testament.

\section{Primary Educational Resources}
No primers or textbooks turned up. Don Burgess's site says its downloads include a Tarahumara/Spanish phrasebook, for non-commercial scholarly and educational use only: \url{https://www.tarahumarabooks.com/}

\section{Linguistic and Reference Materials}
The search found no dictionary or grammar for this variety. The Burgess site above may list some, but I did not check it.

\section{Religious and Community Resources}
\begin{itemize}
\item New Testament ``Onorúgame nila ra'íchali'' (Wycliffe, 2008) on eBible.org: \url{https://ebible.org/details.php?id=tacNT}
\item Scripture Earth listing, with New Testament PDF and app: \url{https://www.scriptureearth.org/00eng.php?iso=tac}
\end{itemize}

\section{Digital and Audio Resources}
\begin{itemize}
\item Complete audio New Testament (MegaVoice): \url{https://megavoice.com/media-cloud/m0b8632-new-testament-tac-western-tarahumara-baja-audio-bible/}
\item Global Recordings audio and slideshows: \url{https://globalrecordings.net/en/language/tac}
\end{itemize}

\section{Curriculum Development Context}
The Scripture text and audio can serve as a corpus for literacy work. Core educational materials such as a dictionary, primer or grammar would have to be sought through further searching or community organizations.
\end{document}